\chapter{Curvature functions on the circle}
\label{chap:curvature}

% archon:covers Gluck/Curvature.lean

The data of the converse theorem is a curvature function preassigned on the
circle $S^1 = \mathbb{R}/2\pi\mathbb{Z}$, modelled as a $2\pi$-periodic function
$\kappa : \mathbb{R} \to \mathbb{R}$.

% SOURCE: [DeTurck-Gluck], §III / §4, p. 9, lines 150-159 (read from references/deturck-gluck-fourvertex.txt)
% SOURCE QUOTE: "FULL CONVERSE TO THE FOUR VERTEX THEOREM. Let κ: S1 → R be a
% continuous function which is either a nonzero constant or else has at least
% two local maxima and two local minima. ... Let κ: S1 → R be any continuous
% strictly positive curvature function"
\begin{definition}\leanok
[Curvature function]
  \label{def:is_curvature_function}
  \lean{Gluck.IsCurvatureFunction}
  \textit{Source: DeTurck--Gluck, §III–4 (strictly positive curvature case).}
  A \emph{curvature function} is a map $\kappa : \mathbb{R} \to \mathbb{R}$ that
  is continuous, $2\pi$-periodic ($\kappa(\theta + 2\pi) = \kappa(\theta)$), and
  strictly positive ($\kappa(\theta) > 0$ for all $\theta$). This is exactly the
  class of preassigned curvatures treated in Gluck's positive case.
\end{definition}

% SOURCE: [DeTurck-Gluck], §4, p. 9, lines 158-174 (read from references/deturck-gluck-fourvertex.txt)
% SOURCE QUOTE: "if s is arc length along this curve, then the equations
% dα/ds = (cos θ, sin θ)  and  dθ/ds = κ(θ)  quickly lead to the explicit formula
% α(θ) = ∫0^θ (cos θ, sin θ) dθ / κ(θ) ."
\begin{definition}\leanok
[Radius of curvature]
  \label{def:radius}
  \lean{Gluck.radius}
  \uses{def:is_curvature_function}
  \textit{Source: DeTurck--Gluck, §4 (the reconstruction integral has integrand $1/\kappa$).}
  The \emph{radius of curvature} of a curvature function $\kappa$ is
  $\rho = 1/\kappa$, i.e. $\rho(\theta) = 1/\kappa(\theta)$. It is the integrand
  weight $1/\kappa(\theta)$ appearing in the source's reconstruction formula
  $\alpha(\theta) = \int_0^\theta e^{i\theta}\,d\theta/\kappa(\theta)$.
\end{definition}

\begin{lemma}\leanok
[Radius is a positive periodic continuous weight]
  \label{lem:radius_props}
  \lean{Gluck.radius_pos}
  \uses{def:radius, def:is_curvature_function}
  If $\kappa$ is a curvature function then $\rho = 1/\kappa$ is continuous,
  $2\pi$-periodic and strictly positive.
\end{lemma}
\begin{proof}

\leanok
  Continuity of $\rho = 1/\kappa$ follows from continuity of $\kappa$ together
  with $\kappa$ being everywhere nonzero (indeed positive). Periodicity is
  inherited from $\kappa$. Positivity of $1/\kappa$ follows from positivity of
  $\kappa$.
\end{proof}

\begin{lemma}

\leanok
[Intermediate value on a closed interval, endpoints in either order]
  \label{lem:ivt_hits}
  \lean{Gluck.ivt_hits}
  A continuous $f : \mathbb{R} \to \mathbb{R}$ on a closed interval $[p, q]$
  (with $p \le q$) attains every value $v$ lying between $f(p)$ and $f(q)$ (in
  either order): if $v \in [\min(f(p), f(q)),\ \max(f(p), f(q))]$ then there is
  $x \in [p, q]$ with $f(x) = v$.
\end{lemma}
\begin{proof}

\leanok
  Project-bespoke packaging of the intermediate value theorem. Split on whether
  $f(p) \le f(q)$ or $f(q) \le f(p)$; in each case rewrite the $\min/\max$ and
  apply the monotone, respectively anti-monotone, form of the intermediate value
  theorem on $[p, q]$.
\end{proof}

% SOURCE: [DeTurck-Gluck], §III, p. 4, lines 13-14 and p. 9 lines 150-153 (read from references/deturck-gluck-fourvertex.txt)
% SOURCE QUOTE: "any continuous realvalued function on the circle which has at
% least two local maxima and two local minima is the curvature function of a
% simple closed curve in the plane."
\begin{definition}\leanok
[Four-vertex condition]
  \label{def:four_vertex_condition}
  \lean{Gluck.FourVertexCondition}
  \uses{def:is_curvature_function}
  \textit{Source: DeTurck--Gluck, Converse statement (§III).}
  A curvature function $\kappa$ satisfies the \emph{four-vertex condition} if
  either
  \begin{itemize}
    \item $\kappa$ is constant; or
    \item $\kappa$ has \emph{two value-separated, alternating local extrema}:
      there are four points
      \[
        p_1 < q_1 < p_2 < q_2, \qquad q_2 < p_1 + 2\pi,
      \]
      lying in a single fundamental period, such that $p_1, p_2$ are local
      maxima of $\kappa$, $q_1, q_2$ are local minima of $\kappa$, and the two
      minimum values lie strictly below the two maximum values:
      \[
        \max\bigl(\kappa(q_1), \kappa(q_2)\bigr) \;<\; \min\bigl(\kappa(p_1), \kappa(p_2)\bigr).
      \]
  \end{itemize}

  \emph{Why the strict value-separation (a project-level correctness refinement
  of the source's phrasing).} The source says ``at least two local maxima and
  two local minima.'' A naive formalization --- ``there exist two distinct points
  that are each a local maximum, and two distinct points that are each a local
  minimum'' using the topological (non-strict) notion of local extremum --- is
  \emph{too weak} and makes the converse theorem false. Counterexample: a
  trapezoidal wave that is constant $=1$ on a bottom plateau, rises to a top
  plateau where it is constant $=2$, and falls back, is unimodal yet has
  infinitely many points that are non-strict local maxima (the whole top
  plateau) and minima (the bottom plateau). It satisfies the naive condition but
  has no four points carrying values $a,b,a,b$ with $a<b$ in cyclic order
  (each level is crossed exactly twice), and indeed is not realizable as a closed
  curve. The genuine content of ``two local maxima and two local minima'' is that
  $\kappa$ oscillates at least twice: two separated humps whose tops strictly
  exceed the two intervening valleys. The condition above encodes exactly that
  and is what the realizability proof consumes (it is precisely equivalent, for
  positive curvature, to the existence of the $a,b,a,b$ interleaving of
  \cref{lem:abab_values}).
\end{definition}

\begin{lemma}\leanok
[Four values $a,b,a,b$ around the circle]
  \label{lem:abab_values}
  \lean{Gluck.exists_abab_of_fourVertex}
  \uses{def:four_vertex_condition, def:is_curvature_function, lem:ivt_hits}
  If a non-constant curvature function $\kappa$ satisfies the four-vertex
  condition, then there exist $0 < a < b$ and four points $\theta_1 < \theta_2 <
  \theta_3 < \theta_4$ lying in a single fundamental period
  ($\theta_4 < \theta_1 + 2\pi$) at which $\kappa$ takes the values
  $a, b, a, b$ in order around the circle.
\end{lemma}
\begin{proof}
  \uses{lem:ivt_hits}
  \leanok
  % SOURCE QUOTE PROOF: "Using the hypothesis that κ has at least two local
  % maxima and two local minima, we find positive numbers 0 < a < b such that
  % κ takes the values a , b , a , b at four points in order around the circle."
  \textit{Source: DeTurck--Gluck, §5, p. 11 (the source asserts this; we prove it
  from the value-separated four-vertex condition).}

  Since $\kappa$ is non-constant, the four-vertex condition gives its second
  disjunct: four points $p_1 < q_1 < p_2 < q_2$ in one period, with $q_2 < p_1 +
  2\pi$, where $p_1, p_2$ are local maxima, $q_1, q_2$ local minima, and
  \[
    \max\bigl(\kappa(q_1), \kappa(q_2)\bigr) \;<\; \min\bigl(\kappa(p_1), \kappa(p_2)\bigr).
  \]

  \emph{Step 1 (the two levels).} Put
  \[
    a := \max\bigl(\kappa(q_1), \kappa(q_2)\bigr), \qquad
    b := \min\bigl(\kappa(p_1), \kappa(p_2)\bigr).
  \]
  The value-separation hypothesis is exactly $a < b$, and $a > 0$ because
  $\kappa > 0$ (so $a \ge \kappa(q_1) > 0$).

  \emph{Step 2 (four sub-arcs, each spanning $[a,b]$).} Consider the four closed
  sub-arcs
  \[
    [p_1, q_1],\quad [q_1, p_2],\quad [p_2, q_2],\quad [q_2, p_1 + 2\pi].
  \]
  Each joins a maximum endpoint (value $\ge b$) to a minimum endpoint (value
  $\le a$); for the last sub-arc the right endpoint value is
  $\kappa(p_1 + 2\pi) = \kappa(p_1) \ge b$ by periodicity. Hence on each sub-arc
  the interval $[a,b]$ lies between the two endpoint values, and by the
  intermediate value theorem (\cref{lem:ivt_hits}) $\kappa$ attains \emph{every}
  value of $[a,b]$ on each sub-arc.

  \emph{Step 3 (select the four points).} Choose
  \[
    \theta_1 \in [p_1, q_1] \text{ with } \kappa(\theta_1) = a, \quad
    \theta_2 \in [q_1, p_2] \text{ with } \kappa(\theta_2) = b,
  \]
  \[
    \theta_3 \in [p_2, q_2] \text{ with } \kappa(\theta_3) = a, \quad
    \theta_4 \in [q_2, p_1 + 2\pi] \text{ with } \kappa(\theta_4) = b,
  \]
  all possible by Step 2. The endpoint inclusions give the weak chain
  \[
    p_1 \le \theta_1 \le q_1 \le \theta_2 \le p_2 \le \theta_3 \le q_2 \le \theta_4 \le p_1 + 2\pi.
  \]
  The assigned values alternate $a, b, a, b$ with $a \ne b$, so consecutive
  chosen points carry different $\kappa$-values, hence are distinct; combined
  with the weak chain this upgrades it to the strict chain $\theta_1 < \theta_2 <
  \theta_3 < \theta_4$.

  \emph{Step 4 (the points lie in one period).} Since $\kappa(p_1) \ge b > a =
  \kappa(\theta_1)$ we have $\theta_1 \ne p_1$, so $\theta_1 > p_1$; together with
  $\theta_4 \le p_1 + 2\pi$ this gives $\theta_4 < \theta_1 + 2\pi$. Thus the four
  points lie in a single fundamental period and $\kappa$ takes the values
  $a, b, a, b$ on them in order around the circle, with $0 < a < b$, as required.
\end{proof}
